\chapter{Neural Networks as Compositional Models}

\section{Layered Architectures}

Neural networks are structured as compositions of simple transformations organized into layers. This layered architecture is the defining feature that enables them to model complex relationships.

\medskip

\noindent
\textbf{Basic structure.}  
A neural network with $L$ layers defines a function
\[
f(x) = f_L \circ f_{L-1} \circ \cdots \circ f_1(x),
\]
where each layer has the form
\[
f_\ell(z) = \sigma(W_\ell z + b_\ell).
\]

\medskip

\noindent
\textbf{Components.}
\begin{itemize}
    \item $W_\ell$: weight matrix
    \item $b_\ell$: bias vector
    \item $\sigma$: nonlinear activation function
\end{itemize}

\medskip

\noindent
\textbf{Forward propagation.}  
Given an input $x$, the network computes intermediate representations:
\[
z^{(1)} = f_1(x), \quad
z^{(2)} = f_2(z^{(1)}), \quad \dots, \quad
z^{(L)} = f(x).
\]

\medskip

\noindent
\textbf{Interpretation.}
\begin{itemize}
    \item Each layer transforms the representation of the data
    \item The final output depends on successive refinements
\end{itemize}

\medskip

\noindent
\textbf{Common activation functions.}
\begin{itemize}
    \item ReLU: $\sigma(z) = \max(0, z)$
    \item Sigmoid: $\sigma(z) = \frac{1}{1 + e^{-z}}$
    \item Tanh: $\sigma(z) = \tanh(z)$
\end{itemize}

\medskip

\noindent
\textbf{Role of nonlinearity.}  
Without nonlinear activation functions, a composition of linear layers reduces to a single linear transformation.

\medskip

\noindent
\textbf{Insight.}  
Nonlinearity enables neural networks to represent complex, non-linear relationships.

\section{Expressivity and Depth}

A central question is how the architecture of a neural network influences its expressive power.

\medskip

\noindent
\textbf{Width vs depth.}
\begin{itemize}
    \item \textbf{Width:} number of units per layer
    \item \textbf{Depth:} number of layers
\end{itemize}

\medskip

\noindent
\textbf{Shallow vs deep networks.}
\begin{itemize}
    \item Shallow networks: few layers, possibly very wide
    \item Deep networks: many layers, often narrower
\end{itemize}

\medskip

\noindent
\textbf{Key observation.}  
Depth allows functions to be represented more efficiently.

\medskip

\noindent
\textbf{Example (compositional functions).}  
Functions with hierarchical structure can be represented compactly by deep networks, but may require exponentially many parameters in shallow networks.

\medskip

\noindent
\textbf{Geometric intuition.}  
Each layer reshapes the input space, gradually transforming complex structures into simpler ones.

\medskip

\noindent
\textbf{Piecewise linearity.}  
Networks with ReLU activations represent piecewise linear functions with potentially many regions.

\medskip

\noindent
\textbf{Insight.}  
Depth increases the number of linear regions exponentially, enhancing expressivity.

\medskip

\noindent
\textbf{Tradeoff.}
\begin{itemize}
    \item Greater depth $\rightarrow$ higher expressivity
    \item But also more challenging optimization
\end{itemize}

\section{Universal Approximation (Insight)}

A fundamental theoretical result explains why neural networks are powerful approximators.

\medskip

\noindent
\textbf{Universal approximation (informal).}  
A feedforward neural network with a single hidden layer and a suitable activation function can approximate any continuous function on a compact domain, given sufficient width.

\medskip

\noindent
\textbf{Interpretation.}
\begin{itemize}
    \item Neural networks are highly expressive
    \item Even shallow networks can approximate complex functions
\end{itemize}

\medskip

\noindent
\textbf{Limitation.}
\begin{itemize}
    \item The theorem does not specify how many units are required
    \item Practical learning may be difficult
\end{itemize}

\medskip

\noindent
\textbf{Role of depth.}  
While shallow networks are theoretically sufficient, deep networks often provide more efficient representations.

\medskip

\noindent
\textbf{Insight.}  
Universal approximation guarantees existence, but not efficiency or learnability.

\medskip

\noindent
\textbf{Connection to practice.}  
Deep architectures are preferred because they:
\begin{itemize}
    \item require fewer parameters,
    \item are easier to train in many cases,
    \item align with compositional structure in data.
\end{itemize}

\section{Representation Learning}

One of the defining features of neural networks is their ability to learn representations automatically.

\medskip

\noindent
\textbf{Learned features.}  
Each layer produces a representation
\[
z^{(\ell)} = f_\ell(z^{(\ell-1)}),
\]
which encodes information about the input.

\medskip

\noindent
\textbf{Hierarchy of representations.}
\begin{itemize}
    \item Early layers: local or low-level features
    \item Intermediate layers: patterns and structures
    \item Final layers: task-specific features
\end{itemize}

\medskip

\noindent
\textbf{Contrast with classical methods.}
\begin{itemize}
    \item Classical models use fixed feature maps
    \item Neural networks learn features from data
\end{itemize}

\medskip

\noindent
\textbf{Geometric viewpoint.}  
Each layer transforms the geometry of the data, making it progressively easier to separate or model.

\medskip

\noindent
\textbf{Compositional learning.}  
Features are built by combining simpler features across layers.

\medskip

\noindent
\textbf{Interpretation.}  
Learning involves discovering both:
\begin{itemize}
    \item how to represent the data,
    \item how to map representations to outputs.
\end{itemize}

\medskip

\noindent
\textbf{Insight.}  
Representation learning is a key reason for the success of deep learning across diverse domains.

\section{A Unifying Perspective}

Neural networks can be understood through three complementary viewpoints:

\begin{itemize}
    \item \textbf{Functional:} compositions of nonlinear functions
    \item \textbf{Geometric:} transformations of data representations
    \item \textbf{Computational:} layered processing architectures
\end{itemize}

\medskip

\noindent
\textbf{Core components.}
\begin{itemize}
    \item Parameters $(W_\ell, b_\ell)$
    \item Nonlinear activations
    \item Compositional structure
\end{itemize}

\medskip

\noindent
\textbf{Key insight.}  
The power of neural networks lies not only in their expressivity, but in their ability to learn structured representations through composition.

\section{Outlook}

This chapter introduced neural networks as compositional models:
\begin{itemize}
    \item layered architectures,
    \item the role of depth in expressivity,
    \item universal approximation as a theoretical foundation,
    \item representation learning as a central principle.
\end{itemize}

\medskip

In the next chapter, we study how neural networks are trained efficiently using backpropagation and automatic differentiation.

\medskip

\noindent
\textbf{Guiding principle.}  
Complex functions can often be understood as compositions of simpler ones, and neural networks provide a framework for learning such compositions from data.